\chapter{Discussion}

This chapter interprets the experimental results presented in
Chapter~4, discusses the limitations of the proposed NLGCN framework,
and positions the work relative to existing methods in the literature.

\section{Interpretation of Results}

\subsection{Why NLGCN Outperforms Baseline Methods}

The experimental results demonstrate that NLGCN consistently achieves
higher Kendall Tau correlation than all three baseline centrality
measures across all four datasets . Several factors
explain this superior performance.

First, traditional centrality measures such as degree centrality and K-Shell decomposition 
rely on a single structural perspective. Degree centrality considers
only immediate connections, while K-Shell assigns a single integer
core number to each node, producing coarse rankings with many tied
scores \cite{chen2012, mukhtar}. NLGCN, by contrast, constructs six
structural channels, three based on Node Local Influence (NLI) and
three based on Node Global Influence (NGI) at multiple aggregation
orders, that jointly encode both local and global structural
properties of each node. This richer
representation allows the model to distinguish nodes that simple
centrality measures would rank equally.

Second, the multi-scale aggregation strategy captures hierarchical
structural information that single-hop methods miss \cite{zhang2019gcn}. By constructing $W_{\text{NLI}1}$, $W_{\text{NLI}2}$,
$W_{\text{NLI}3}$ and their NGI counterparts through iterative
neighbourhood summation, the model encodes not only a node's own
structural position but also the cumulative influence of its
two-hop neighbourhood \cite{guo2026}. This is particularly beneficial
in sparse networks such as Budapest (avg.\ degree 4.17), where local
degree information alone is insufficient to distinguish influential
nodes, explaining NLGCN's largest margin of improvement
(+46.0\% over K-Shell) on that dataset.

\subsection{Role of the Channel Attention Mechanism}

The channel attention mechanism plays a critical
role in NLGCN's performance by learning to weight the six structural
channels dynamically. Not all channels contribute equally to influence
prediction across different network types. In dense networks such as
C.~elegans (avg.\ degree 14.46), higher-order NGI channels that
capture long-range structural influence are likely upweighted, as
global position matters more in well-connected graphs. In sparse
networks like Budapest, local NLI channels may receive higher
attention weights, since immediate neighbourhood structure is more
discriminative when global paths are long.

This adaptive weighting makes NLGCN more robust than fixed-weight
approaches. Without the attention
mechanism, all six channels would contribute equally to the
convolution, reducing the model's ability to specialise its
representation for different network structures.

\section{Limitations}

Despite its strong empirical performance, the proposed NLGCN
framework has several limitations that should be acknowledged.

\subsection{Scalability}

The neighbourhood matrix construction requires selecting the top-$L$
neighbours for each node and building a fixed-size subgraph
representation \cite{guo2026, tu2025}. For the datasets used in this
work (297--834 nodes), this is computationally feasible. However,
for very large networks with tens of thousands or millions of nodes,
this process becomes a bottleneck, as it requires computing all-pairs
shortest path distances for NGI (complexity $O(n^3)$ with
Floyd-Warshall or $O(nm)$ with Brandes' algorithm
\cite{farooq2018, chen2012}). Extending NLGCN to large-scale networks
would require approximations such as landmark-based distance
estimation or sampling-based neighbourhood construction.

\subsection{Sensitivity to Hyperparameters}

The model's performance depends on the choice of $L$ (neighbourhood
size) and $\alpha$ (balancing constant in NGI) \cite{guo2026}.
A larger $L$ captures more structural context but increases the
tensor size and computational cost. A poorly chosen $\alpha$ may
over- or under-weight the degree contribution in the NGI formula.
A grid search over these parameters, as performed in the original
NLGCN paper \cite{guo2026}, would likely yield further performance
improvements.


\subsection{No Ablation Experiments}

The current implementation does not include ablation experiments
to isolate the contribution of individual components
(attention mechanism, NLI channels, NGI channels, multi-scale
aggregation). While the full model performs strongly, it is not
possible to determine from the current results how much each
component contributes independently \cite{guo2026, tu2025}.

\subsection{Static Graph Assumption}

The NLGCN framework assumes a static, unweighted, undirected graph.
Real-world networks such as social networks and transportation
systems are dynamic, nodes and edges are added or removed over
time \cite{hafiene2020}. The current model does not account for
temporal evolution of network structure, which may limit its
applicability to time-varying influence prediction tasks.

\section{Comparison with Related Work}

\subsection{Heuristic Centrality Methods}

Traditional methods for influential node detection including
degree centrality, betweenness centrality,
closeness centrality , and K-Shell decomposition are computationally efficient and interpretable,
but have well-documented limitations \cite{chen2012, mukhtar}.
They rely on a single structural dimension and fail to capture the
multi-scale patterns that determine real spreading dynamics
\cite{guo2026, thotakura2024}. The Mukhtar et al.\ \cite{mukhtar}
analysis using ANOVA shows that no single centrality measure
consistently outperforms others across different network types,
motivating the need for combined or learning-based approaches.
Chen et al.\ \cite{chen2012} demonstrated that semi-local
centrality provides a better trade-off between computational
complexity and ranking accuracy than global measures, but still
falls short of learning-based methods on complex networks.

\subsection{GNN-Based Approaches}

Graph Neural Network-based methods \cite{zhang2019gcn}
have advanced the state of the art in influence prediction by
learning node representations directly from graph structure.
InfGCN \cite{tu2025} and related models integrate centrality
features as node attributes and process them through GCN layers,
achieving improved ranking accuracy over heuristic baselines.
Tu et al.\ \cite{tu2025} proposed CGCAN, which combines CNN and
GNN modules with an attention mechanism, demonstrating that
hybrid architectures that jointly capture local and global topology
outperform single-path networks. Patel and Singh \cite{patel2025}
proposed SE\_GCN, which uses structural equivalence and degree
sequences to construct feature matrices processed by GCN layers,
achieving strong Kendall Tau correlations on synthetic and
real-world networks.

NLGCN \cite{guo2026} is most directly related to these hybrid
approaches. Its key distinction is the explicit construction of
multi-scale NLI and NGI features embedded into neighbourhood
matrices as image-like channels, enabling a standard 2D CNN
to extract structural patterns without requiring a full GCN
message-passing stack \cite{guo2026, zhang2019gcn}. This makes
NLGCN computationally lighter at inference time compared to
full GCN-based approaches, while still achieving competitive
ranking accuracy.

\subsection{Survey Perspective}

Hafiene et al.\ \cite{hafiene2020} provide a comprehensive survey
of influential node detection methods in both static and dynamic
networks, categorizing approaches into heuristic, greedy, and
hybrid methods. Their analysis highlights that static methods
applied to dynamic networks lose performance over time, and that
the most accurate approaches combine structural features with
simulation-based ground truth,exactly the paradigm adopted
by NLGCN.